% =====================================================================
%  五、模型建立与求解（核心章节）
%  每个子问题统一采用「建模思路 → 模型建立 → 算法/求解 → 结果分析」结构。
% =====================================================================
\section{模型建立与求解}

% ---------------------------------------------------------------------
\subsection{问题一：示向度交会定位模型}
\subsubsection{建模思路}
\textbf{【待填充：说明由两条含误差的示向度交会出定位区域的基本思想。】}

\subsubsection{模型建立}
设检测点 $P_1(x_1,y_1)$、$P_2(x_2,y_2)$ 处测得示向度分别为 $\theta_1$、$\theta_2$，
考虑测量误差 $\delta$，则真实方位落在 $[\theta_i-\delta,\theta_i+\delta]$ 内，两条
扇形射线束的交叠区域即为干扰源的可能位置区域。
\textbf{【待填充：写出射线/扇形边界的解析表达式，定义交会定位区域。】}

区域\textbf{直径}定义为区域内任意两点间距离的上确界，等价于凸包的直径，可用
旋转卡壳（Rotating Calipers）算法在 $O(n\log n)$ 内求得；“直径圆”能否覆盖区域
则归结为\textbf{最小覆盖圆}（smallest enclosing circle）问题。
\textbf{【待填充：给出直径与最小覆盖圆的数学定义及判据。】}

\subsubsection{求解与结果}
\textbf{【待填充：给出问题一的求解流程、示例交会区域示意图（图~2）与结论。】}

% ---------------------------------------------------------------------
\subsection{问题二：第二检测点最优选择模型}
\subsubsection{建模思路}
\textbf{【待填充：说明“距离近”与“交会角接近 $90^\circ$”之间的权衡。】}

\subsubsection{模型建立}
在距离 $d$ 处，示向度误差 $\delta$ 产生的横向定位误差约为 $d\tan\delta$。两条
示向度交会时，可用\textbf{Fisher 信息矩阵}刻画定位不确定度，其行列式
$\det(\mathbf{F})$ 与交会角及两检测点到目标的距离有关；等价地可用 GDOP 或
误差椭圆面积作为准则。候选区域可由等精度线 / 扇环刻画。
\textbf{【待填充：写出 Fisher 信息矩阵 / 误差椭圆的具体形式与最优化目标。】}

\subsubsection{求解与结果}
\textbf{【待填充：给出第二检测点选择策略、候选区域图示（图~3）与结果。】}

% ---------------------------------------------------------------------
\subsection{问题三：多源快速定位与清除策略}
\subsubsection{建模思路}
\textbf{【待填充：说明“先粗扫建频道台账，再精定位收敛频道”的信息驱动策略，以及
时间预算的核算（移动、测向、换频道、清除各自耗时）。】}

\subsubsection{模型建立}
将测向机视为在 $1800\,\mathrm{m}$ 半径圆域内移动的智能体，待清除的 $10$--$16$
个干扰源位置未知。该问题可建模为\textbf{信息驱动的旅行商问题（TSP with active
sensing）}：以栅格或正六边形铺排搜索路线，用射线束交点与置信椭圆修剪候选区域，
依据“频道—射线束”台账决定下一次测向位置与频道。
\textbf{【待填充：给出目标函数（如总清除时间最小 / 清除率最大）、状态转移与
决策规则。】}

\subsubsection{求解与结果}
\textbf{【待填充：给出算法流程（可配伪代码）、模拟器联调结果、清除比例与平均
定位清除时间等实测指标，附表格（表~2）与轨迹图（图~4）。】}

% ---------------------------------------------------------------------
\subsection{问题四：定向干扰源定位模型}
\subsubsection{建模思路}
\textbf{【待填充：说明定向源方向未知带来的“测不到信号”多重成因与诊断逻辑。】}

\subsubsection{模型建立}
定向干扰源仅在有限方位覆盖角内可被检测，故“未测得信号”需区分三种原因：该频道
无干扰源、超出有效接收半径、处于定向覆盖角之外。建模需在问题三的搜索框架上引入
\textbf{定向覆盖角与源朝向}两个未知量，并设计相应的绕测 / 诊断策略。
\textbf{【待填充：给出定向源的可检测条件、参数估计模型与诊断规则。】}

\subsubsection{求解与结果}
\textbf{【待填充：给出问题四的求解流程与结果，并与问题三对比。】}
